\paragraph{Correction to the pseudoisotopy input.}
The split formula stated in Hatcher's Theorem 3.1 must not be used
in that generality.  The corrected smooth sequence, for a compact
smooth $M$ with $\dim M\geq5$ and $\pi_2M=0$, is
\[
 \operatorname{Wh}_3(\pi_1M)
 \xrightarrow{\chi_{\rm Wh}}
 \operatorname{Wh}_1^+(\pi_1M;\mathbb Z/2)
 \longrightarrow\pi_0P^{\rm Diff}(M,\partial M)
 \longrightarrow\operatorname{Wh}_2(\pi_1M)
 \longrightarrow0.
\]
Here $P^{\rm Diff}(M,\partial M)$ consists of diffeomorphisms of
$M\times I$ fixed on $M\times\{0\}\cup\partial M\times I$,
with the usual $C^\infty$ topology.  Igusa's \emph{Generalized
Grassmann invariant---redrawn}, arXiv:2502.19147v2,
Theorem 7.1, printed p.23, states this sequence; Remark 7.3, p.26,
explicitly warns against the asserted splitting even when
$\chi_{\rm Wh}=0$.  The introduction, p.2, identifies
\[
 \operatorname{Wh}_3(\pi)=
 K_3(\mathbb Z[\pi])/
 \bigl(K_3(\mathbb Z)+\operatorname{im}
             [\Omega^{\rm fr}_3(B\pi)\to K_3(\mathbb Z[\pi])]\bigr).
\]

For $\pi=\mathbb Z^n$, this correction does not alter the torus
calculation.  Iterating the Laurent fundamental theorem for the
regular ring $\mathbb Z$ gives, with $t_1,\dots,t_n$ the coordinate
units,
\[
\begin{split}
 K_3(\mathbb Z[\mathbb Z^n])\cong {}&K_3(\mathbb Z)
 \oplus\bigoplus_i K_2(\mathbb Z)[t_i]
 \oplus\bigoplus_{i<j} K_1(\mathbb Z)[t_i][t_j]\\
 &\oplus\bigoplus_{i<j<k}K_0(\mathbb Z)[t_i][t_j][t_k].
\end{split}
\]
The splitting maps are products with the displayed unit classes.
The unit of algebraic $K$-theory sends the framed class
$\eta\in\pi^s_1$ to $[-1]\in K_1(\mathbb Z)$, and sends
$\eta^2$ to $\{-1,-1\}$, which generates
$K_2(\mathbb Z)=\mathbb Z/2$.  The unit in degree zero generates
$K_0(\mathbb Z)=\mathbb Z$.  Accordingly every nonconstant summand
in the last display belongs to the framed assembly image:
use $\eta^2$ times a coordinate circle, $\eta$ times a coordinate
2-torus, and the fundamental framed 3-tori, respectively.  The
constant summand is killed separately in Igusa's definition.
Thus $\operatorname{Wh}_3(\mathbb Z^n)=0$.  The same Laurent
calculation in degree two, with the monomial/framed generators,
gives $\operatorname{Wh}_2(\mathbb Z^n)=0$.  There are no Nil
summands because all the rings are regular.  Finally the conjugation
action of $\mathbb Z^n$ is trivial, and hence
\[
 \pi_0P^{\rm Diff}(T^n)
 \cong\operatorname{Wh}_1^+(\mathbb Z^n;\mathbb Z/2)
 = (\mathbb Z/2)[\mathbb Z^n]/(\mathbb Z/2)[0]
 =\bigoplus_{a\in\mathbb Z^n\setminus\{0\}}\mathbb Z/2.
\]
This uses an exact sequence with both outside groups zero, not a
splitting claim.  For the present application one can take $n\geq6$:
Courte--Porcelli, arXiv:2506.06110, Section 6.2, p.34,
equation (6.12) and Lemmas 6.16--6.17, independently use this torus
calculation in that range.  Their footnote on p.4 records that the
$n=5$ mapping-class statement in Hatcher's Theorem 4.1 was attributed
to announced work, so the larger range avoids that separate issue.

\paragraph{The unstable category comparison.}
Igusa's theorem just quoted is a smooth theorem.  The required
comparison with PL is an independent, unstabilized result of
Burghelea--Lashof, \emph{The homotopy type of the space of
diffeomorphisms II}, Trans.\ AMS 196 (1974), Section 6,
Theorem 6.1, p.41, proof pp.42--43.  For a compact connected smooth
manifold $M$, with the underlying PL structure supplied by a smooth
triangulation, the comparison induces
\[
 \pi_0P^{\rm Diff}(M,\partial M)
       \xrightarrow{\ \sim\ }
 \pi_0P^{\rm PL}(M,\partial M).
\]
We use it only in dimension at least six.  In their proof the
relative smoothing space is a section space over $M$ with fibre
\[
 E(\mathrm{PL}_{n+1}/O_{n+1},\mathrm{PL}_n/O_n),
\]
the space of paths from the chosen base point to the indicated
subspace.  Their Theorem 5.1, p.37, and Lemma 6.4(c), p.42,
give vanishing of its section-space $\pi_0$ and $\pi_1$;
the relative homotopy sequence therefore gives the displayed
component isomorphism without stabilizing $M$.
Theorem 6.2, p.41, proof p.43, gives
\[
 \pi_iP^{\rm PL}(M,\partial M)
       \xrightarrow{\ \sim\ }
 \pi_iP^{\rm TOP}(M,\partial M)\qquad(i\geq0,\ n>5).
\]
Consequently the computed direct sum is also the actual PL and
TOP component group of $T^n$, $n\geq6$.
Hatcher's discussion after Theorem 3.1, p.7, is consistent with this
geometric comparison, but no general split exact sequence from that
older theorem is needed.  Courte--Porcelli Lemma 6.16 concerns
stable spaces, whereas Lemma 6.17 concerns the smooth stabilization;
neither alone is being used as an unstabilized PL comparison.
